\documentclass{article}
\usepackage{amsmath}   % equations
\usepackage{amssymb}   % \mathbb, also loads amsfonts
\usepackage{graphicx}  % images
\usepackage{array}     % for \raggedright\arraybackslash in tabular
\usepackage{booktabs}
\usepackage{cleveref}
\title{Feynmann Excercises Physics
}
% \author{nicolas.kubail}
% \date{August 2025}
\pdfminorversion=7
\pdfobjcompresslevel=0
\begin{document}

\maketitle


\section{Atoms in Motion}
Refer to chapters 1-3.

\subsection{If heat is merely molecular motion, what is the difference between a hot, stationary baseball and a cool, rapidly moving one?}
\begin{itemize}
    \item A hot stationary baseball is made of atoms that are moving and bouncing around at a high speed, even though the macro object of the “baseball” is standing still.
    \item The cool rapidly moving one is the opposite, it’s constituent atoms are moving slowly, while the macro object is moving fast across space.
    \item The difference is in the manner of the motion, organized (thrown ball), disorganized (heat).
\end{itemize}

Every atom's velocity can be split into two independent pieces:

  $$\vec{v}{\text{atom}} = \vec{v}{\text{bulk}} + \vec{v}_{\text{random}}$$
\begin{itemize}
    \item $\vec{v}_{\text{bulk}}$ — the shared, coherent velocity every atom has in addition to its jiggling, because the whole object is moving together (the
    thrown ball case).
    \item $\vec{v}_{\text{random}}$ — the individual, chaotic jiggling on top of that, different in direction and magnitude for every atom (the heat).
\end{itemize}

\subsection{If the atoms of all objects are perpetually in motion, how can there be any permanent objects, such as fossil imprints?}
\begin{itemize}
    \item $E_{\text{bond}}$ = \textbf{bond energy}, this is the energy needed to completely separate two bonded atoms. pull them from their equilibrium spacing out to infinity, overcoming the attractive force between them. It's a fixed property of the pair of atoms and the type of bond (ionic, covalent, etc.) — it doesn't depend on temperature.
    \item $E_{\text{thermal}}$ = \textbf{thermal energy} is the typical amount of kinetic energy an atom has just from jiggling around due to heat, at tempearature $T$. The standard estimate (average) is: $$E_{\text{thermal}} \sim k_B T$$
    where $k_B$ is Boltzmann's constant, and $T$ is the absolute temperature (Kelvin). 
    \item Whether an atom stays bonded or breaks free comes down to comparing these two:
    \subitem If $E_{\text{thermal}} \ll E_{\text{bond}}$ (i.e. $k_B T \ll E_{\text{bond}}$): the atom doesn't have enough energy to escape the attractive well. It
    vibrates in place but stays bonded — this is the solid, "permanent" case.
    \subitem If $E_{\text{thermal}} \gtrsim E_{\text{bond}}$: the atom's jiggling is energetic enough to break free of the bond — this is melting, evaporation, or
    bond-breaking.

    \item At everyday/room temperature, for most solids, $k_B T$ is far smaller than typical bond energies, which is exactly why ordinary objects (and fossil imprints) stay intact indefinitely without special intervention and why it takes a large energy input (heat, chemical reaction, or impact) to disrupt that arrangement.
\end{itemize}

\subsection{Explain qualitatively why and how friction in a moving machine produces heat. Explain also, if you can, why heat cannot produce useful motion by the reverse process.}

\begin{itemize}
    \item Heat is nothing more than the chaotic movement of atoms $\vec{v}_{\text{random}}$, the more movement the more heat. 
    \item When two solids are rubbing against each other, the (organized) mechanical energy going into the system becomes (disorganized) internal energy in the material and the surrounding environment.
    \item The microscopic imperfections of the solids are grinding against each other. So the builk kinetic energy gets converted into heat at the atomic level.
    \item The second part of the question has to do with entropy (second law of thermodynamics). A system can go from an ordered state to a disordered state, but it is highly improbably for spontaneous organization to happen in order to reverse that process. We usually need an external mechanism to go from disordered state to useful motion.
\end{itemize}

\subsection{Chemists have found that the molecules of rubber consist of long criss-crossed chains of atoms. Explain why a rubber band becomes warm when it is stretched.}

\begin{itemize}
    \item High entropy at rest: long polymer chains are tangled (high configurational entropy, lots of possible arrangements).
    \item Low entropy when deformed: chains are pulled taut and thus ordered (lower configurational entropyfewer possible arrangements, lower disorder).
    \item The loss of entropy needs to be conserved, so it appears inside the materail as extra atomic jiggling.
    \item More atomic jiggling is heat, so the band is warmer.
\end{itemize}

\subsection{What should happen to a rubber band which is supporting a given weight, if it is heated? (To find out, try it.)}

\begin{itemize}
    \item You add heat $Q>0$.
    \item This forces the molecules to move toward a higher entropy state $S>0$.
    \item The higher entropy state for a polymer is tangled and coiled.
    \item Coiled chains shorten the band, causing it to contract and lift the weight.
\end{itemize}

\subsection{Can you explain why there are no crystals that have the shape of a regular pentagon? (Triangles, squares, and hexagons are common in crystal forms.)}

\begin{itemize}
    \item Because you can’t pack pentagons on the Euclidean plane, as 2pi is not divisible by 108 degrees. So you can’t tesselate a plane. (No translational symmetry atomic structure.
\end{itemize}

\subsection{You are given a large number of steel balls of equal diameter d and a container of known volume V. Every dimension of the container is much greater than the diameter of a ball. What is the greatest number of balls, N, that can be placed in the container}

\begin{itemize}
    \item $4/3 pi r^3 = S$ (volume of sphere)
    \item However packing a 3d space with spheres will create gaps
    \item So what is the tesselation pattern of perfect spheres? We need to calculate the in between space 
    \item N = V/(S)
    \item N = V/(S+X)
    \item Gauss’s sphere packing rule
\end{itemize}

\subsection{How should the pressure P of a gas vary with n, the number of atoms per unit volume, and (v), the average speed of an atom? (Should P be proportional to n and/ or to (v), or should it vary more, or less, rapidly than linearly?)}

\begin{itemize}
    \item P = pressure
    \item n = atoms/unit volume
    \item v = average speed
    \item average speed is heat so the more it heats up the greater the pressure regardless of density (n)
    \item P also goes up with more density regardless of velocity staying constant
    \item $P = 1/3 N/V mv^2$
    \item Linearly with more density, and quadratically with average speed

\end{itemize}


\section{Placholder}


Placholder text
$$P^* = \arg\min_{P} \sum_{i=1}^{n-1} w(v_i, v_{i+1})$$


\end{document}
